\documentclass[14pt]{extarticle}
\usepackage{fontspec}
\usepackage{polyglossia}
\usepackage{amsmath}
\usepackage{amssymb}
\usepackage{xcolor}
\usepackage{fancyhdr}
\usepackage{graphicx}
\usepackage{listings}
\usepackage[most]{tcolorbox}
% \usepackage{setspace}
\usepackage{pifont}
\usepackage{enumitem}
\usepackage{titlesec}
\usepackage[bottom]{footmisc}
\usepackage{titling}
\usepackage{minted}
\usepackage{etoolbox}
\usepackage{array}
\usepackage{longtable} % page-breakable tables (must load before polyglossia/bidi RTL setup)
\usepackage{extsizes}
\usepackage{float}
\usepackage{booktabs}
\usepackage{hyperref}
\usepackage[onehalfspacing]{setspace}
\newif\ifwhiteblack
\ifdefined\setwhiteblack
  \setwhiteblack
\fi

\ifwhiteblack
  \usepackage[a4paper, landscape, top=1.25cm, bottom=1.25cm, left=1.25cm, right=1.25cm, headheight=1.5cm, headsep=0.5cm]{geometry}
\else
  \usepackage[a4paper, top=2.7cm, bottom=2.7cm, left=2.5cm, right=2.5cm, headheight=1.5cm, headsep=0.5cm]{geometry}
\fi
\usepackage{tikz}
\usepackage{datetime2}

\usetikzlibrary{automata, positioning, arrows.meta, arrows}
\providecommand{\LTcaptype}{table}

\newfontfamily\emoji{Segoe UI Emoji}

\pagestyle{fancy}

\setmainlanguage[numerals=western]{arabic}
\setotherlanguage{english}
\newfontfamily\arabicfont[Script=Arabic]{Calibri}
\newfontfamily\arabicfonttt[Script=Arabic]{Courier New}



\ifwhiteblack
\lstset{
  language=[Sharp]C,
  numbers=left,
  stepnumber=1,
  numbersep=8pt,
  frame=single,
  basicstyle=\ttfamily\small,
  keywordstyle=\bfseries,
  stringstyle=\ttfamily,
  commentstyle=\itshape
}
\else
\lstset{
  language=[Sharp]C,
  numbers=left,
  stepnumber=1,
  numbersep=8pt,
  frame=single,
  basicstyle=\ttfamily\small,
  keywordstyle=\color{blue},
  stringstyle=\color{red},
  commentstyle=\color{green!50!black}
}
\fi

\newif\ifdetailed
\ifdefined\setdetailed
  \setdetailed
\fi

\newif\ifwithsols
\ifdefined\setwithsols
  \setwithsols
\fi

\ifwhiteblack
  \AtBeginDocument{
    \LARGE
    \usemintedstyle{bw}
  }
\fi

% Default minted look: compact code that wraps long lines.
% Not applied in black-and-white landscape mode, which already uses \LARGE.
\ifwhiteblack\else
  \setminted{fontsize=\small, breaklines=true}
\fi

% unified tcolorboxes for articles
\ifwhiteblack
\tcbset{colback=white, colframe=black, fonttitle=\bfseries, boxrule=1pt}
\newtcolorbox{boxDef}[1][]{title={{\emoji📘} تعريف\ifx\\#1\\\else ~#1\fi :}}
\newtcolorbox{boxExercise}[1][]{title={{\emoji🧩} تمرين\ifx\\#1\\\else ~#1\fi :}}
\newtcolorbox{boxExample}[1][]{title={{\emoji📝} مثال\ifx\\#1\\\else ~#1\fi :}}
\newtcolorbox{boxNote}[1][]{title={{\emoji✨} ملاحظة\ifx\\#1\\\else ~#1\fi :}}
\newtcolorbox{boxAttention}[1][]{title={{\emoji🔔} تنبيه\ifx\\#1\\\else ~#1\fi :}}
\newtcolorbox{boxWarning}[1][]{title={{\emoji⚡} ملاحظة هامة\ifx\\#1\\\else ~#1\fi :}}
\newtcolorbox{boxSolution}[1][]{title={{\emoji✅} حل\ifx\\#1\\\else ~#1\fi :}}
\newtcolorbox{boxSymbol}[1][]{title={{\emoji🔣} رمز\ifx\\#1\\\else ~#1\fi :}}
\newtcolorbox{boxHint}[1][]{title={{\emoji💡} تلميح\ifx\\#1\\\else ~#1\fi :}}
\else
\tcbset{colback=white, colframe=black, fonttitle=\bfseries, boxrule=0.8pt}
\newtcolorbox{boxDef}[1][]{colback=blue!5!white,colframe=blue!75!black,
  title={{\emoji📘} تعريف\ifx\\#1\\\else ~#1\fi :}}
\newtcolorbox{boxExercise}[1][]{colback=cyan!5!white,colframe=cyan!70!black,
  title={{\emoji🧩} تمرين\ifx\\#1\\\else ~#1\fi :}}
\newtcolorbox{boxExample}[1][]{colback=yellow!5!white,colframe=orange!90!black,
  title={{\emoji📝} مثال\ifx\\#1\\\else ~#1\fi :}}
\newtcolorbox{boxNote}[1][]{colback=gray!10!white,colframe=black,
  title={{\emoji✨} ملاحظة\ifx\\#1\\\else ~#1\fi :}}
\newtcolorbox{boxAttention}[1][]{colback=magenta!10!white,colframe=magenta!80!black,
  title={{\emoji🔔} تنبيه\ifx\\#1\\\else ~#1\fi :}}
\newtcolorbox{boxWarning}[1][]{colback=red!5!white,colframe=red!75!black,
  title={{\emoji⚡} ملاحظة هامة\ifx\\#1\\\else ~#1\fi :}}
\newtcolorbox{boxSolution}[1][]{colback=green!5!white,colframe=green!60!black,
  title={{\emoji✅} حل\ifx\\#1\\\else ~#1\fi :}}
\newtcolorbox{boxSymbol}[1][]{colback=purple!5!white,colframe=purple!70!black,
  title={{\emoji🔣} رمز\ifx\\#1\\\else ~#1\fi :}}
\newtcolorbox{boxHint}[1][]{colback=teal!5!white,colframe=teal!60!black,
  title={{\emoji💡} تلميح\ifx\\#1\\\else ~#1\fi :}}
\fi


\ifwhiteblack
\tcbset{simplecode/.style={ colback=white, colframe=black, boxrule=1pt, arc=2pt, left=4pt,right=4pt,top=4pt,bottom=4pt}}
\else
\tcbset{simplecode/.style={ colback=gray!5, colframe=black!50, boxrule=0.4pt, arc=2pt, left=4pt,right=4pt,top=4pt,bottom=4pt}}
\fi
% Compact code boxes: tighter line spacing (the preamble uses onehalfspacing) and small vertical gaps.
% Must come AFTER the simplecode definition above.
\tcbset{simplecode/.append style={fontupper=\setstretch{0.95}, top=4pt, bottom=4pt, before skip=6pt, after skip=6pt}}
\newenvironment{boxCode}{\begin{tcolorbox}[simplecode]}{\end{tcolorbox}}

% Code cell used inside tables (keeps code left-to-right inside RTL tables).
% Use inside a fixed-width p{...} column: \begin{pycell}\begin{minted}...\end{minted}\end{pycell}
\newtcolorbox{pycell}{frame empty,colback=white,boxsep=0pt,left=2pt,right=2pt,top=2pt,bottom=2pt,width=\linewidth}

\newcolumntype{C}[1]{>{\centering\arraybackslash}p{#1}}

\makeatletter
\DTMnewdatestyle{dotted}{%
  \renewcommand{\DTMdisplaydate}[4]{\two@digits{##3}.\two@digits{##2}.##1}%
  \renewcommand{\DTMDisplaydate}{\DTMdisplaydate}%
}
\makeatother
\DTMsetdatestyle{dotted}

% redefine spaces after titles
\makeatletter
\renewcommand{\@maketitle}{%
  \begin{center}
    {\huge \bfseries \@title \par}%
    \vskip 0.2em % space between title and author
    {\large \@author \par}%
    % \vskip 0.2em % space between author and date
    % {\normalsize \@date \par}%
    {\small \hfill \begin{english}\DTMtoday\end{english} \par}%
  \end{center}
}
\makeatother

\fancyhf{} % clear default
\fancypagestyle{plain}{
  \fancyhf{}
  \fancyhead[L]{مدرسة التسامح الشاملة}
  % \fancyhead[L]{\includegraphics[height=1cm]{../../../images/logoTasamoh.png}}
  \fancyhead[R]{الأستاذ محمود اغبارية}
  \fancyfoot[C]{\thepage}
}

\fancyhead[L]{مدرسة التسامح الشاملة}
\fancyhead[R]{الأستاذ محمود اغبارية}
\fancyfoot[C]{\thepage}

\setcounter{tocdepth}{3} % only section subsection and subsubsection in TOC

\newcommand{\importpdfpage}[2]{%
    \thispagestyle{empty}
    \newgeometry{left=0cm,right=0cm,top=0cm,bottom=0cm}
    \includegraphics[page=#2,width=0.95\paperwidth,height=0.95\paperheight,keepaspectratio]{#1}
    \restoregeometry
}

\newcommand{\insertFullImg}[1]{%
    \noindent
    \makebox[\textwidth][c]{
        \includegraphics[width=0.9\paperwidth,keepaspectratio]{#1}%
    }%
}%


\setminted{fontsize=\small, baselinestretch=0.95, breaklines=true}

\title{كراس تمارين متدرج: الوراثة وتعدد الأشكال}

\begin{document}

\maketitle
\thispagestyle{fancy}

\begin{boxDef}[مشروع نظام الموظفين في شركة "السعادة للجميع"]
في شركة \textbf{"السعادة للجميع"}، ترغب الإدارة في بناء نظام برمجي متكامل لإدارة شؤون الموظفين، حساب رواتبهم، وتنظيم عمليات التصويت في الشركة. \\
في هذا الكراس، سنبني النظام \textbf{خطوة بخطوة}: سنبدأ ببناء فئات منفصلة بدون وراثة، ثم نكتشف التكرار ونعالجه عبر بناء فئة أساسية وتطبيق \textbf{الوراثة}، وبعدها نوسّع النظام ليشمل إدارة الشركة وتطبيق \textbf{تعدد الأشكال} و\textbf{إعادة تعريف العمليات}.
\end{boxDef}

\begin{boxWarning}
\textbf{تعليمات هامة للعمل:}
في \textbf{كل تمرين} تغيّر فيه شيئاً في الفئات أو تضيف عملية جديدة، \textbf{يجب عليك} تحديث البرنامج الرئيسي \textenglish{Main} لإنشاء كائنات واستدعاء العمليات الجديدة وفحص صحّة عملها على الشاشة.
\end{boxWarning}

% ==============================================================
\section*{المرحلة الأولى: بناء الفئات (مقدمة للوراثة)}
% ==============================================================

\begin{enumerate}[label=\textbf{تمرين \arabic*:}, wide=0pt, leftmargin=*, itemsep=2.2em]

    \item \textbf{بناء فئة تقني الحواسيب (\textenglish{Technician}) بدون وراثة} \\
    ابنِ فئة باسم \textenglish{Technician} لإدارة بيانات تقني الحواسيب في الشركة. \\
    \textbf{الصفات (\textenglish{private}):}
    \begin{itemize}
        \item \textenglish{id} (\textenglish{string}): رقم هوية الموظف.
        \item \textenglish{name} (\textenglish{string}): اسم الموظف.
        \item \textenglish{seniority} (\textenglish{int}): أقدمية الموظف في الشركة (عدد السنوات).
        \item \textenglish{computers} (\textenglish{int}): عدد الحواسيب التي تقع تحت مسؤوليته.
    \end{itemize}
    \textbf{العمليات (\textenglish{public}):}
    \begin{itemize}
        \item \textbf{عملية بنائية أولى:} تتلقى رقم الهوية، الاسم، سنوات الأقدمية، وعدد الحواسيب، وتهيئ جميع الصفات.
        \item \textbf{عملية بنائية ثانية (لتقني جديد):} تتلقى رقم الهوية، الاسم، وعدد الحواسيب فقط، وتهيئ الأقدمية (\textenglish{seniority}) إلى \textenglish{0}.
        \item \textbf{عمليات \textenglish{Get}:} لجميع الصفات (\textenglish{GetId}، \textenglish{GetName}، \textenglish{GetSeniority}، \textenglish{GetComputers}).
        \item \textbf{عمليات \textenglish{Set}:}
        \begin{itemize}
            \item \textenglish{SetName(string name)}: لتعديل اسم الموظف.
            \item \textenglish{SetComputers(int computers)}: لتعديل عدد الحواسيب (بشرط أن يكون العدد \textbf{غير سالب} $\ge 0$؛ إذا كانت القيمة سالبة لا يتم التعديل).
        \end{itemize}
        \item \textbf{العملية \textenglish{AddSeniorityYear()}:} تزيد عدد سنوات أقدمية الموظف بسنة واحدة ($+1$).
        \item \textbf{العملية \textenglish{ToString()}:} تعيد نصاً يصف بيانات التقني.
    \end{itemize}
    \textbf{الفحص في \textenglish{Main}:} \\
    اكتب في \textenglish{Main} مقطع برنامج ينشئ كائنين من الفئة \textenglish{Technician} (أحدهما بالعملية البنائية الأولى والآخر بالثانية)، يستدعي \textenglish{AddSeniorityYear()} و \textenglish{SetComputers}، ويطبع وصفهما باستخدام \textenglish{ToString()}.

    \item \textbf{بناء فئة عضو لجنة العمال (\textenglish{UnionMember}) بدون وراثة} \\
    ابنِ فئة باسم \textenglish{UnionMember} لإدارة بيانات موظف يشغل عضوية في لجنة عمال الشركة. \\
    \textbf{الصفات (\textenglish{private}):}
    \begin{itemize}
        \item \textenglish{id} (\textenglish{string}): رقم هوية الموظف.
        \item \textenglish{name} (\textenglish{string}): اسم الموظف.
        \item \textenglish{seniority} (\textenglish{int}): أقدمية الموظف في الشركة (عدد السنوات).
        \item \textenglish{unionYears} (\textenglish{int}): عدد سنوات عضويته في لجنة العمال.
    \end{itemize}
    \textbf{العمليات (\textenglish{public}):}
    \begin{itemize}
        \item \textbf{عملية بنائية:} تتلقى رقم الهوية، الاسم، سنوات الأقدمية، وسنوات العضوية في اللجنة، وتهيئ جميع الصفات.
        \item \textbf{عمليات \textenglish{Get}:} لجميع الصفات (\textenglish{GetId}، \textenglish{GetName}، \textenglish{GetSeniority}، \textenglish{GetUnionYears}).
        \item \textbf{عمليات \textenglish{Set}:}
        \begin{itemize}
            \item \textenglish{SetName(string name)}: لتعديل اسم الموظف.
            \item \textenglish{SetUnionYears(int unionYears)}: لتعديل سنوات العضوية في اللجنة (بشرط أن تكون القيمة \textbf{غير سالبة} ولا تتجاوز أقدميته في الشركة: $0 \le \text{unionYears} \le \text{seniority}$).
        \end{itemize}
        \item \textbf{العملية \textenglish{AddSeniorityYear()}:} تزيد عدد سنوات أقدمية الموظف في الشركة بسنة واحدة ($+1$).
        \item \textbf{العملية \textenglish{ToString()}:} تعيد نصاً يصف بيانات عضو اللجنة.
    \end{itemize}
    \textbf{الفحص في \textenglish{Main}:} \\
    أضف إلى \textenglish{Main} إنشاء كائن من الفئة \textenglish{UnionMember}، وجرّب استدعاء عملياته (\textenglish{AddSeniorityYear}، \textenglish{SetUnionYears}، \textenglish{ToString}) وطباعة النتائج.

    \item \textbf{استخلاص المشترك وبناء فئة الموظف الأساسي (\textenglish{Employee})} \\
    راجع الفئتين اللتين بنيتهما في التمرين 1 والتمرين 2. جد جميع الصفات والعمليات \textbf{المتطابقة (المشتركة فقط)} بين الفئتين، وأنشئ فئة جديدة باسم \textenglish{Employee} تتضمن فقط هذه الصفات والعمليات المشتركة. \\
    \textbf{الفحص في \textenglish{Main}:} \\
    أنشئ كائناً من الفئة \textenglish{Employee} في \textenglish{Main}، وافحص عملياته واطبع بياناته.

\end{enumerate}

\begin{boxAttention}[- نقطة للتفكير]
هل كان من الأريح برمجياً تطوير الجزء المشترك مرة واحدة فقط، ثم إضافة الجزء الفريد لكل فئة على حدة؟ \\
ربما سيساعدنا ذلك أيضاً عند بناء فئات لأنماط إضافية من الموظفين (مثل مسؤول مشاريع، محاسب، مدير فرع...). \\
\textbf{والآن ننتقل إلى تعلم الوراثة...}
\end{boxAttention}

\clearpage

% ==============================================================
\section*{المرحلة الثانية: تطبيق الوراثة بين الفئات}
% ==============================================================

\begin{enumerate}[resume, label=\textbf{تمرين \arabic*:}, wide=0pt, leftmargin=*, itemsep=2em]

    \item \textbf{بناء علاقة الوراثة بين \textenglish{Employee} و \textenglish{Technician} و \textenglish{UnionMember}}
    \begin{itemize}
        \item \textbf{أولاً:} حدّث الفئة \textenglish{Employee} وفق قواعد البرمجة كائنية التوجه بحيث تشكّل \textbf{فئة أساسية (فئة أب)} للفئات التي سترث منها (\textenglish{Technician} و \textenglish{UnionMember})، وذلك بتحويل أذونات الوصول لصفاتها إلى محمية (\textenglish{protected}).
        \item \textbf{ثانياً:} أعد بناء الفئتين \textenglish{Technician} و \textenglish{UnionMember} بحيث ترثان من الفئة \textenglish{Employee} (باستخدام \textenglish{: Employee}):
        \begin{itemize}
            \item أبقِ في كل فئة فرعية الصفات والعمليات الخاصة بها \textbf{فقط} (واحذف التكرار الموجود في \textenglish{Employee}).
            \item استعن بالعملية البنائية لفئة الأب باستخدام الكلمة المحجوزة \textenglish{base}.
            \item في العملية \textenglish{ToString()} لكل فئة فرعية، استدعِ \textenglish{base.ToString()} وألحق بها الصفة الخاصة بالفئة الفرعية.
        \end{itemize}
    \end{itemize}

    فيما يلي ملخص واجهة الفئة الأساسية \textenglish{Employee} بعد التحديث:
    \begin{center}
    \renewcommand{\arraystretch}{1.2}
    {\small
    \begin{tabular}{|p{0.44\textwidth}|p{0.48\textwidth}|}
    \hline
    \multicolumn{2}{|c|}{\textbf{الفئة الأساسية: \textenglish{Employee}}} \\
    \hline
    صفات محمية: رقم الهوية، الاسم، وسنوات الأقدمية. & \textenglish{protected string id, name;} \newline \textenglish{protected int seniority;} \\
    \hline
    عمليتان بنائيتان لتهيئة صفات الموظف. & \textenglish{public Employee(string id, string name, int seniority)} \newline \textenglish{public Employee(string id, string name)} \\
    \hline
    عمليات \textenglish{Get}، \textenglish{SetName}، \textenglish{AddSeniorityYear}، و \textenglish{ToString}. & \textenglish{GetId(), GetName(), GetSeniority(), SetName(...), AddSeniorityYear(), ToString()} \\
    \hline
    \end{tabular}
    }
    \end{center}

    وملخص واجهتي الفئتين الفرعيتين الوارثتين \textenglish{Technician} و \textenglish{UnionMember}:
    \begin{center}
    \renewcommand{\arraystretch}{1.2}
    {\small
    \begin{tabular}{|p{0.44\textwidth}|p{0.48\textwidth}|}
    \hline
    \multicolumn{2}{|c|}{\textbf{الفئة الفرعية: \textenglish{Technician : Employee}}} \\
    \hline
    صفة خاصة لعدد الحواسيب. & \textenglish{private int computers;} \\
    \hline
    عمليتان بنائيتان تستعينان بـ \textenglish{base}. & \textenglish{public Technician(string id, string name, int seniority, int computers)} \newline \textenglish{public Technician(string id, string name, int computers)} \\
    \hline
    عمليات الحواسيب وإعادة تعريف \textenglish{ToString}. & \textenglish{GetComputers(), SetComputers(int c), override ToString()} \\
    \hline
    \multicolumn{2}{|c|}{\textbf{الفئة الفرعية: \textenglish{UnionMember : Employee}}} \\
    \hline
    صفة خاصة لسنوات العضوية في اللجنة. & \textenglish{private int unionYears;} \\
    \hline
    عملية بنائية تستعين بـ \textenglish{base}. & \textenglish{public UnionMember(string id, string name, int seniority, int unionYears)} \\
    \hline
    عمليات العضوية وإعادة تعريف \textenglish{ToString}. & \textenglish{GetUnionYears(), SetUnionYears(int u), override ToString()} \\
    \hline
    \end{tabular}
    }
    \end{center}

    \textbf{الفحص في \textenglish{Main}:} \\
    افحص الفئات المحدّثة في \textenglish{Main}: أنشئ كائنات من الفئات الثلاث، واستدعِ على كائني \textenglish{Technician} و \textenglish{UnionMember} عمليات موروثة من \textenglish{Employee} (مثل \textenglish{GetName} و \textenglish{AddSeniorityYear}) وعمليات خاصة بكل منهما، واطبع الكائنات.

    \item \textbf{إضافة فئة فرعية ثالثة: مسؤول المشاريع (\textenglish{Supervisor})} \\
    أضف فئة جديدة باسم \textenglish{Supervisor} تمثل مسؤول مشاريع في الشركة يدير فريقاً مكوّناً من 15 موظفاً كحد أقصى، وترث من الفئة \textenglish{Employee}. \\
    \textbf{الصفات الخاصة الإضافية (\textenglish{private}) في \textenglish{Supervisor}:}
    \begin{itemize}
        \item \textenglish{team} (\textenglish{Employee[]}): مصفوفة تمثل فريق الموظفين الذين يشرف عليهم (بحجم أقصى 15).
        \item \textenglish{teamCount} (\textenglish{int}): عدد الموظفين الفعليين حالياً في فريقه.
    \end{itemize}
    \textbf{العمليات (\textenglish{public}):}
    \begin{itemize}
        \item \textbf{عملية بنائية أولى:} تتلقى رقم الهوية، الاسم، وسنوات الأقدمية، وتهيئ صفات الأب عبر \textenglish{base}، وتنشئ مصفوفة \textenglish{team} بحجم 15، وتهيئ \textenglish{teamCount} إلى \textenglish{0}.
        \item \textbf{عملية بنائية ثانية:} تتلقى رقم الهوية، الاسم، سنوات الأقدمية، ومصفوفة موظفين جاهزة \textenglish{Employee[] team}، وتهيئ صفات الأب عبر \textenglish{base}، وتنسخ الحدود غير الفارغة (\textenglish{!= null}) إلى مصفوفة الفريق وتحدّث \textenglish{teamCount}.
        \item \textbf{عمليات \textenglish{Get}:} للصفتين (\textenglish{GetTeam()} و \textenglish{GetTeamCount()}).
        \item \textbf{العملية \textenglish{AddEmployeeToTeam(Employee emp)}:} تتلقى كائن موظف (من نمط \textenglish{Employee} أو أي فئة ترث منه)، وإذا كان هناك متسع في فريق المشروع ($\text{teamCount} < 15$) تضيفه إلى الفريق وتعيد \textenglish{true}، وإلا تعيد \textenglish{false}.
        \item \textbf{العملية \textenglish{TransferBonusSeniority(Employee emp)}:} عملية تشجيعية تتلقى كائن موظف \textenglish{emp} وتستدعي عليه \textenglish{AddSeniorityYear()} لمنحه سنة أقدمية إضافية، بشرط أن يكون هذا الموظف عضواً موجوداً فعلياً داخل مصفوفة \textenglish{team} الخاصة بمسؤول المشاريع الحالي (تعيد \textenglish{true} إذا وجدته في فريقه ومنحته السنة، و \textenglish{false} خلاف ذلك).
        \item \textbf{العملية \textenglish{ToString()}:} تعيد وصفاً نصياً لمسؤول المشاريع بالاستعانة بـ \textenglish{base.ToString()} مع إضافة عدد أعضاء فريقه الحاليين.
    \end{itemize}
    \textbf{الفحص في \textenglish{Main}:} \\
    أنشئ كائناً من الفئة \textenglish{Supervisor} في \textenglish{Main}، وأضف إلى فريقه موظفاً عادياً (\textenglish{Employee})، وتقنياً (\textenglish{Technician})، وعضو لجنة (\textenglish{UnionMember})، ثم جرّب العملية \textenglish{TransferBonusSeniority} واطبع بيانات المسؤول وأعضاء فريقه.

\end{enumerate}

\clearpage

% ==============================================================
\section*{المرحلة الثالثة: تعدد الأشكال الأساسي (التحويل لأعلى)}
% ==============================================================

\begin{enumerate}[resume, label=\textbf{تمرين \arabic*:}, wide=0pt, leftmargin=*, itemsep=2em]

    \item \textbf{بناء فئة إدارة الشركة (\textenglish{Company})} \\
    أضف فئة جديدة لمشروعنا باسم \textenglish{Company}. تمثّل هذه الفئة مجموعة موظفي الشركة من جميع الأنماط، بالاعتماد على الفئة الأساسية \textenglish{Employee} (أي تشمل أيضاً كل كائن من فئة ترث من \textenglish{Employee}).
    \begin{itemize}
        \item \textbf{الصفات:}
        \begin{itemize}
            \item \textenglish{public const int MAX = 100}: ثابت يمثل الحد الأقصى للموظفين في الشركة.
            \item \textenglish{private Employee[] employees}: مصفوفة لحفظ موظفي الشركة.
            \item \textenglish{private int current}: عدد الموظفين الفعليين حالياً في الشركة.
        \end{itemize}
        \item \textbf{عملية بنائية \textenglish{public Company()}:} تنشئ المصفوفة بحجم \textenglish{MAX} وتهيئ \textenglish{current} إلى \textenglish{0}.
        \item \textbf{عملية \textenglish{GetCurrent()}:} تعيد عدد الموظفين الحاليين في الشركة.
    \end{itemize}
    \textbf{الفحص في \textenglish{Main}:} \\
    أنشئ كائناً من الفئة \textenglish{Company} في \textenglish{Main} واطبع عدد الموظفين الابتدائي فيه.

    \item \textbf{إضافة موظف إلى الشركة (\textenglish{AddEmployee})} \\
    اكتب في الفئة \textenglish{Company} عملية باسم \textenglish{AddEmployee(Employee e)} تتلقى موظفاً من أي نمط وتضيفه إلى مصفوفة موظفي الشركة: افحص أولاً إن كان هناك متسع في المصفوفة؛ إذا نعم، أضف الموظف وأعد \textenglish{true}، وإلا فلا تضفه وأعد \textenglish{false}. \\
    \textbf{الفحص في \textenglish{Main}:} \\
    أضف في \textenglish{Main} عدة موظفين من الأنماط الأربعة (\textenglish{Employee}، \textenglish{Technician}، \textenglish{UnionMember}، و \textenglish{Supervisor}) إلى كائن الشركة \textenglish{Company}.

    \item \textbf{البحث عن بيانات موظف حسب رقم الهوية (\textenglish{GetEmployeeDetails})} \\
    اكتب في الفئة \textenglish{Company} عملية باسم \textenglish{GetEmployeeDetails(string id)} تتلقى رقم هوية موظف (افترض أنه فريد) وتعيد سلسلة نصية تصف بيانات هذا الموظف (استخدم \textenglish{ToString()}). إذا لم يوجد موظف بهذا الرقم في الشركة، تعيد العملية نصاً فارغاً \textenglish{""}. \\
    \textbf{الفحص في \textenglish{Main}:} \\
    استدعِ العملية في \textenglish{Main} لرقم هوية موجود وآخر غير موجود واطبع النتيجة.

    \item \textbf{عدّ الموظفين الذين يملكون نفس الأقدمية (\textenglish{CountSameSeniority})} \\
    اكتب في الفئة \textenglish{Company} عملية باسم \textenglish{CountSameSeniority(string id)} تتلقى رقم هوية موظف (افترض أنه فريد وموجود في الشركة)، وتعيد عدد الموظفين في الشركة الذين لديهم نفس عدد سنوات الأقدمية التي يمتلكها صاحب هذا الرقم. \\
    \textbf{الفحص في \textenglish{Main}:} \\
    استدعِ العملية في \textenglish{Main} واطبع النتيجة.

    \item \textbf{تصفية الموظفين حسب الأقدمية (\textenglish{GetEmployeesByMinSeniority})} \\
    اكتب في الفئة \textenglish{Company} عملية باسم \textenglish{GetEmployeesByMinSeniority(int minYears)} تتلقى حداً أدنى لسنوات الأقدمية، وتعيد مصفوفة من نمط \textenglish{Employee[]} تحتوي على جميع موظفي الشركة الذين أقدميتهم أكبر أو تساوي \textenglish{minYears} (بحجم مطابق تماماً لعددهم). إذا لم يوجد أي موظف يحقق الشرط، يُعاد \textenglish{null}. \\
    \textbf{الفحص في \textenglish{Main}:} \\
    استدعِ العملية في \textenglish{Main} واطبع بيانات جميع الموظفين في المصفوفة المُعادة.

    \item \textbf{إيجاد الموظف صاحب الأقدمية الأكبر (\textenglish{GetMostVeteranId})} \\
    اكتب في الفئة \textenglish{Company} عملية باسم \textenglish{GetMostVeteranId()} تعيد رقم هوية (\textenglish{id}) الموظف صاحب أكبر أقدمية في الشركة (افترض وجود موظف واحد كهذا، وإذا كانت الشركة خالية تعيد \textenglish{""}). \\
    \textbf{الفحص في \textenglish{Main}:} \\
    استدعِ العملية في \textenglish{Main} واطبع رقم هوية الموظف الأقدم في الشركة.

\end{enumerate}

\clearpage

% ==============================================================
\section*{المرحلة الرابعة: إعادة تعريف العمليات (\textenglish{virtual} و \textenglish{override})}
% ==============================================================

\begin{boxNote}[- قواعد حساب الرواتب ونقاط التصويت واستحقاق المكافأة]
تُتخذ القرارات في الشركة عبر التصويت باستخدام \textbf{نقاط التأثير}، كما يتحدد الراتب الشهري والمكافأة السنوية لكل موظف بحسب نمطه وأقدميته كما يلي:
\begin{center}
\renewcommand{\arraystretch}{1.4}
{\small
\begin{tabular}{|p{0.20\textwidth}|p{0.37\textwidth}|p{0.35\textwidth}|}
\hline
\textbf{نمط الموظف} & \textbf{الراتب الشهري (\textenglish{GetSalary})} & \textbf{نقاط التصويت (\textenglish{GetScore})} \\
\hline
\textbf{موظف عادي} \newline (\textenglish{Employee}) &
راتب أساسي $5000$ ش.ج + $500$ ش.ج عن كل سنة أقدمية: \newline
$5000 + 500 \cdot \text{seniority}$ &
$4$ نقاط أساسية + نقطة عن كل سنة أقدمية: \newline
$4 + \text{seniority}$ \\
\hline
\textbf{تقني حواسيب} \newline (\textenglish{Technician}) &
راتب الموظف العادي + $250$ ش.ج عن كل 10 حواسيب كاملة: \newline
$\text{base} + \lfloor \frac{\text{computers}}{10} \rfloor \cdot 250$ &
نقاط الموظف العادي + نقطة عن كل حاسوب: \newline
$\text{base} + \text{computers}$ \\
\hline
\textbf{عضو لجنة عمال} \newline (\textenglish{UnionMember}) &
يماثل راتب الموظف العادي تماماً. &
\textbf{ضعف} نقاط الموظف العادي + نقطتان عن كل سنة باللجنة: \newline
$2 \cdot \text{base} + 2 \cdot \text{unionYears}$ \\
\hline
\textbf{مسؤول مشاريع} \newline (\textenglish{Supervisor}) &
راتب الموظف العادي + $100$ ش.ج عن كل موظف فعلي في فريقه: \newline
$\text{base} + 100 \cdot \text{teamCount}$ &
نقاط الموظف العادي + مجموع نقاط كافة أعضاء فريقه الفعليين. \\
\hline
\end{tabular}
}
\end{center}
\end{boxNote}

\begin{enumerate}[resume, label=\textbf{تمرين \arabic*:}, wide=0pt, leftmargin=*, itemsep=2em]

    \item \textbf{تطبيق العمليات القابلة لإعادة التعريف: \textenglish{GetSalary()} و \textenglish{GetScore()}}
    \begin{enumerate}[label=\textbf{(\arabic*)}]
        \item أضف إلى الفئة الأساسية \textenglish{Employee} العمليتين العامتين القابلتين لإعادة التعريف:
        \begin{itemize}
            \item \textenglish{public virtual double GetSalary()}
            \item \textenglish{public virtual int GetScore()}
        \end{itemize}
        وأعد تعريفهما (\textenglish{override}) في الفئات الفرعية \textenglish{Technician} و \textenglish{UnionMember} و \textenglish{Supervisor} وفق الجدول أعلاه، مع الحرص على استدعاء \textenglish{base.GetSalary()} و \textenglish{base.GetScore()} لتجنب تكرار المعادلات الحسابية.
        \item أضف إلى الفئة \textenglish{Company} العملية \textenglish{public double TotalSalary()} التي تحسب وتعيد مجموع الرواتب الشهرية لجميع موظفي الشركة.
    \end{enumerate}
    \textbf{الفحص في \textenglish{Main}:} \\
    في \textenglish{Main}، اطبع راتب ونقاط تصويت كل موظف، ثم اطبع مجموع الرواتب الكلي في الشركة عبر \textenglish{TotalSalary()}.

    \item \textbf{إضافة عملية استحقاق المكافأة السنوية (\textenglish{YearlyBonus})}
    \begin{enumerate}[label=\textbf{(\arabic*)}]
        \item أضف لمشروع الشركة عملية عامة قابلة لإعادة التعريف تعريفها: \textenglish{public virtual bool YearlyBonus()}
        \begin{itemize}
            \item تُعرّف العملية لأول مرة في الفئة \textenglish{Employee} وتعيد \textenglish{true} إذا كانت أقدمية الموظف 5 سنوات على الأقل ($\text{seniority} \ge 5$)، وإلا تعيد \textenglish{false}.
            \item في الفئة \textenglish{Technician} يُعاد تعريف العملية (\textenglish{override}) وتعيد \textenglish{true} إذا كان التقني يحقق شرط الأقدمية (عبر \textenglish{base.YearlyBonus()}) \textbf{و}يتولى مسؤولية 100 حاسوب على الأقل ($\text{computers} \ge 100$)، وإلا تعيد \textenglish{false}.
            \item في الفئتين \textenglish{UnionMember} و \textenglish{Supervisor} يُعاد تعريف العملية (\textenglish{override}) بحيث تعيد \textenglish{false} دائماً (لأن أعضاء اللجنة ومسؤولي المشاريع غير مشمولين في هذه المكافأة).
        \end{itemize}
        \item أضف إلى الفئة \textenglish{Company} عملية باسم \textenglish{GetBonusEmployees()} تعيد مصفوفة (\textenglish{Employee[]}) تضم جميع موظفي الشركة المستحقين للمكافأة السنوية (بحسب نتيجة العملية \textenglish{YearlyBonus()}).
    \end{enumerate}
    \textbf{الفحص في \textenglish{Main}:} \\
    استدعِ في \textenglish{Main} العملية \textenglish{GetBonusEmployees()} واطبع بيانات الموظفين المستحقين للمكافأة.

\end{enumerate}

\clearpage

% ==============================================================
\section*{المرحلة الخامسة: تعدد الأشكال المتقدم (\textenglish{is} والتحويل لأسفل)}
% ==============================================================

\begin{enumerate}[resume, label=\textbf{تمرين \arabic*:}, wide=0pt, leftmargin=*, itemsep=2em]

    \item \textbf{تخصيص علاوة مسؤول المشاريع وتصفية التقنيين (\textenglish{GetSeniorTechnicians})}
    \begin{enumerate}[label=\textbf{(\arabic*)}]
        \item \textbf{تحديث \textenglish{GetSalary()} في الفئة \textenglish{Supervisor}:} \\
        قررت إدارة الشركة تعديل منحة الإدارة الخاصة بمسؤول المشاريع بحيث يحصل على راتب الموظف العادي (\textenglish{base.GetSalary()}) + علاوة قدرها \textbf{$100$ ش.ج عن كل تقني حواسيب (\textenglish{Technician})} ضمن فريقه + علاوة قدرها \textbf{$50$ ش.ج عن كل عضو لجنة عمال (\textenglish{UnionMember})} ضمن فريقه (دون علاوة عن الموظفين العاديين في الفريق). حدّث العملية \textenglish{GetSalary()} في \textenglish{Supervisor} باستخدام العامل \textenglish{is}.
        \item \textbf{عملية \textenglish{GetSeniorTechnicians} في الفئة \textenglish{Company}:} \\
        اكتب عملية في الفئة \textenglish{Company} تعريفها:
        \begin{center}
        \textenglish{public Technician[] GetSeniorTechnicians(int minComputers)}
        \end{center}
        تعيد العملية مصفوفة من نمط \textenglish{Technician[]} تحتوي فقط على تقنيي الحواسيب في الشركة الذين عدد الحواسيب تحت مسؤوليتهم أكبر أو يساوي \textenglish{minComputers} (بهدف ترقيتهم مستقبلاً).
    \end{enumerate}
    \textbf{الفحص في \textenglish{Main}:} \\
    اطبع راتب مسؤول المشاريع في \textenglish{Main} بعد التعديل، ثم استدعِ \textenglish{GetSeniorTechnicians} واطبع بيانات التقنيين الذين أرجعتهم العملية.

    \item \textbf{تقرير المراجعة الإدارية (\textenglish{AuditReport}) وعملية التصويت العامة (\textenglish{IsAccepted})}
    \begin{enumerate}[label=\textbf{(\arabic*)}]
        \item \textbf{العملية \textenglish{AuditReport} في الفئة \textenglish{Company}:} \\
        اكتب عملية لقسم الموارد البشرية في الفئة \textenglish{Company} باسم \textenglish{public string[] AuditReport()}. تعيد العملية مصفوفة نصوص (\textenglish{string[]})، بحيث يصف كل سطر فيها موظفاً في الشركة من نمط \textbf{تقني حواسيب (\textenglish{Technician}) أو عضو لجنة عمال (\textenglish{UnionMember})} ممن يُصنَّفون كـ \textbf{"كادر تشغيلي متمرس"} بالمعطيات التالية:
        \begin{center}
        \textenglish{\$"Type:\{empType\}, ID:\{id\}, Name:\{name\}, Score:\{score\}"}
        \end{center}
        \begin{itemize}
            \item \textenglish{empType} هو نمط الموظف: \textenglish{"Technician"} أو \textenglish{"Union"}.
            \item \textenglish{id} هو رقم هوية الموظف، و \textenglish{name} هو اسمه، و \textenglish{score} هو نقاط تصويته.
            \item في حالة كان الموظف تقني حواسيب (\textenglish{Technician})، يجب إضافة عدد حواسيبه في نهاية النص: \textenglish{", Computers:\{computers\}"}.
            \item في حالة كان الموظف عضو لجنة (\textenglish{UnionMember})، يجب إضافة سنوات عضويته باللجنة في نهاية النص: \textenglish{", UnionYears:\{unionYears\}"}.
        \end{itemize}
        \textbf{انتبه:} الموظف الذي يُصنَّف كـ \textbf{"كادر تشغيلي متمرس"} يجب أن تكون أقدميته في الشركة أكثر من 3 سنوات ($\text{seniority} > 3$)، وأن تكون نقاط تصويته (\textenglish{GetScore()}) أكبر من 15 نقطة!

        \item \textbf{العملية العامة \textenglish{IsAccepted} في فئة \textenglish{Program}:} \\
        طُرح اقتراح رسمي في الشركة لتغيير اسمها من "السعادة للجميع" إلى \textbf{"السعادة الآن وفي المستقبل"}. \\
        صوّت جميع الموظفين العاديين والتقنيين \textbf{تأييداً} للاقتراح (مع)، بينما عارض جميع أعضاء لجنة العمال ومسؤولي المشاريع هذا الاقتراح (\textbf{ضد}). \\
        يُعتمد القرار فقط إذا كان مجموع نقاط المؤيدين \textbf{أكبر تماماً} من مجموع نقاط المعارضين. \\
        اكتب في فئة \textenglish{Program} العملية الثابتة التي تعريفها:
        \begin{boxCode}
        \begin{english}
        \begin{minted}{csharp}
public static bool IsAccepted(object[] arr)
        \end{minted}
        \end{english}
        \end{boxCode}
        تتلقى العملية مصفوفة كائنات عامة \textenglish{object[]} تمثل موظفي الشركة، وتعيد \textenglish{true} إذا قُبل الاقتراح، وإلا تعيد \textenglish{false}.
    \end{enumerate}
    \textbf{الفحص في \textenglish{Main}:} \\
    في \textenglish{Main}، استدعِ \textenglish{AuditReport()} واطبع جميع أسطر التقرير، ثم أنشئ مصفوفة \textenglish{object[]} تضم موظفي الشركة ومررها إلى \textenglish{IsAccepted} واطبع نتيجة التصويت.

\end{enumerate}

\vspace{1cm}
\begin{flushleft}
أرجو لكم وقتًا ممتعًا.

الأستاذ محمود اغبارية.
\end{flushleft}

\end{document}
